\sscp{\tsc{West \ili{Central} Flores}}{03-02-03}
The members of \tsc{Flores} (exclusive of \tsc{Flores-Lembata},
see \srf{03-02-04}) can be placed in a single node on the basis
of lowering of high vowels to mid before final *R in a subset of words.
Examples of the lowering of high vowels before final *R are given
in \trf{03-28}, along with cognates in other members of \tsc{\ili{Bima}-Lembata} to
show that this change does not occur in these nodes.

\begin{table}\tcp{\tsc{West \ili{Central} Flores} *-uR > \it{o}; **-iR > \it{e}}{03-28}
	\begin{threeparttable}\st{5pt}
	\begin{tabular}{lllllr@{ }l}\lsptoprule
*gloss	&	egg	&	tail	&	coconut	&	grain-head	& 	\mc{2}{l}{water}\\
%			&		&		&			&			&	*wahiyəR > \\
PMP									&	*qatəl\tbr{uR}			&	*ik\tbr{uR}			&	*ni\tbr{uR}					&	*bul\tbr{iR}		&*wahiyəR >	&	**wah\tbr{iR}\\ \midrule
\ili{Komodo}							&	{\hqa}təl\tbr{o}		&	{\hr}ik\tbr{o}	&	{\hr}ni\tbr{u}\su{‡}&									&&	\cg{(banu)}\\
\ili{Manggarai}						&	{\hqa}təl\tbr{o}		&	{\hr}ik\tbr{o}	&	{\hr}ni\tbr{o}			&	\cg{(kenar)}		&&	{\hr}ʋa\tbr{e}\\
Kepo{\Q}						&	{\hq}\cg{(ruɣa)}		&	{\hr}ik\tbr{o}	&	{\hr}ni\tbr{o}			&									&&	{\hr}wa\tbr{e}ʔ\\
\ili{Waerana}							&	{\hqa}təl\tbr{o}ʔ		&	{\hr}ik\tbr{o}ʔ	&	{\hr}ni\tbr{o}ʔ			&	{\hr}wol\tbr{e}ʔ&&	{\hr}wa\tbr{e}ʔ\\
\ili{Rembong}							&	{\hqa}təl\tbr{o}ʔ		&	{\hr}ik\tbr{o}ʔ	&	{\hr}ni\tbr{o}ʔ			&	\cg{(təŋuʔ)}		&&	{\hr}wa\tbr{e}ʔ\\
\ili{Namut}								&	\hp{qa}‹tel\tbr{o}›	&	‹ək\tbr{o}›			&											&									&&	{\hr}wa\tbr{e}\\
Palu{\Q}e						&	{\hqa}təl\tbr{o}		&	{\hr}iʔ\tbr{o}	&	{\hr}ni\tbr{o}			&									&&	{\hr}va\tbr{e}\\
\ili{Lio}									&	{\hqa}təl\tbr{o}		&	{\hr}ek\tbr{o}	&	{\hr}ni\tbr{o}			&	{\hr}ʋol\tbr{e}	&&	{\hr}a\tbr{e}\\
\ili{Ende}								&	{\hqa}təɹ\tbr{o}		&	{\hr}ʔek\tbr{o}	&	{\hr}ni\tbr{o}			&	{\hr}woɹ\tbr{e}	&&	{\hr}ʔa\tbr{e}\\
C. \ili{Nage}							&	{\hqa}təl\tbr{o}		&	{\hr}ek\tbr{o}	&	{\hr}ni\tbr{o}			&									&&	{\hr}a\tbr{e}\\
Nga{\Q}o, \ili{Oja}				&	{\hqa}təd͡ʎ̝\tbr{o}		&	{\hr}ek\tbr{o}	&	{\hr}ni\tbr{o}			&									&&	{\hr}a\tbr{e}\\
%Nga{\Q}o, Watumite	&	{\hqa}təd͡ʎ̝\tbr{o}		&	{\hr}ek\tbr{o}	&	{\hr}ni\tbr{o}			&									&&	{\hr}a\tbr{e}\\
East \ili{Nage}						&	{\hqa}təd\tbr{o}		&	{\hr}ek\tbr{o}	&	{\hr}ni\tbr{o}			&									&&	{\hr}a\tbr{e}\\
\ili{Keo}, Ud.						&	{\hqa}təd\tbr{o}		&	{\hr}ʔek\tbr{o}	&	{\hr}ni\tbr{o}			&	{\hr}wod\tbr{e}	&&	{\hr}ʔa\tbr{e}\\
\ili{Ngadha}							&	{\hqa}təl\tbr{o}		&	{\hr}ek\tbr{o}	&	{\hr}ni\tbr{o}			&	{\hr}vol\tbr{e}	&&	{\hr}va\tbr{e}\\
\ili{Rongga}							&	{\hqa}təl\tbr{o}		&	{\hr}ek\tbr{o}	&	{\hr}ni\tbr{o}			&	{\hr}wol\tbr{e}	&&	{\hr}wa\tbr{e}\\ \midrule
\ili{Bima}								&	{\hqa}dolu					&	\cg{(keto)}			&	{\hr}niʔu						&	{\hr}wuri				&&	{\hr}oi\\
\ili{Sika}								&	{\hqa}təlo\su{†}		&	{\hr}iʔur				&	{\hr}niur						&	{\hr}βulir			&&	{\hr}βair\\
\ili{Kambera}							&	{\hqa}tilu					&	{\hr}kiku				&	\cg{(kokur)}				&	{\hr}wili				&&	{\hr}wɐi\\
		\lspbottomrule
	\end{tabular}
		\begin{tablenotes}\tnsize

			\item[†]	{Final *R > {\0} and *u > \it{o} indicates that \ili{Sika} \it{təlo} ‘egg’ is a loan.}
			\item[‡]	{\ili{Komodo} \it{niu} ‘coconut’ may be a loan from \ili{Bima} \it{niʔu}
								and/or final vowel lowering may be blocked due to the previous high vowel.}
		\end{tablenotes}
	\end{threeparttable}
\end{table}

The main exception to this change is that it does not always
occur when the penultimate vowel is also high.
One example is *bibiR > \ili{Manggarai} \it{ʋiʋir},
Nga{\Q}o \it{ʋiʋi} ‘lips’, among other reflexes.
There are also cases without a penultimate high vowel in which
such lowering does not occur.
One example is *qauR ‘bamboo’ > \ili{Manggarai} \it{aur},
\ili{Lio} \it{au} ‘bamboo’, among other reflexes.
While this sound change sets \tsc{West \ili{Central} Flores} apart
from other \tsc{\ili{Bima}-Lembata} languages,
the same sound change is found in the \tsc{Rote-Meto} languages
of \tsc{Timor-Babar} (\crf{ch:TB}), as well as the \tsc{\ili{Central} Timor} languages (\crf{ch:CT}).

In between \tsc{West Flores}
and \tsc{\ili{Central} Flores} are several poorly described languages
including: \ili{Manus} (possibly a variety of \ili{Manggarai}),
\ili{Rajong}, and \ili{Namut}.
In the absence of sufficient data which is both accessible
and reliable, we cannot determine whether these languages belong
with \tsc{West Flores} or \tsc{\ili{Central} Flores}.
We place them for the time being as isolates within \tsc{West \ili{Central} Flores}.
Although not the focus of his work,
\citet[2--5]{Asplund2020} provides the most detailed summary of these
languages, including some remarks on the possible status
and relationships of these languages.
